%% introduction.tex
%%

%% ==============================
\section{Introduction}
\label{sec:Introduction}
\subsection{Background and Motivation}
As the demand for compute becomes larger than ever, data centers become a very critical part of modern digital infrastructure, because it gives organizations access to a scale of computational power which, in most cases, cannot be managed on-premises. To manage these physical machines means to efficiently assign arriving requests or tasks to Virtual Machines, while in accordance with SLA violations and maximizing hardware utilization. The problem now comes down to finding the best mode of allocation in a heterogeneous cloud environment.

With the emergence of AI/ML workloads, it becomes more important than ever to find a way to efficiently and consistently assign workloads to the right Virtual Machine Instance. These days, data centers are beginning to face heavier batch jobs that are also compute-heavy. Modern data center managers are also responsible for balancing trade-offs that aim to reduce makespan and financial costs, while also trying to maintain high hardware utilization \cite{Jeba2019Green}.

\subsection{Problem Statement}
Traditional scheduling algorithms such as Heterogeneous Earliest Finish Time, Min-Min and Max-Min have fixed rules that they follow, in terms of task metadata and resource availability; their strengths become their shortcomings in the context of a heterogeneous cloud environment. Meta-heuristic algorithms like Simulated Annealing use higher-level workflows to explore complex search spaces \cite{KUMAR20191}, but the problem comes from their high computational complexity which makes them unsuitable for high-velocity workloads.

As mentioned above, static or traditional scheduling algorithms cannot simply adapt to constantly changing workloads, and when this happens, resources are underutilized, SLAs are violated,  and operational costs balloon

\subsection{Proposed Solution}
This project presents a dynamic Multi-Armed Bandit (MAB) hyper-heuristic framework for task scheduling heterogeneous cloud environments. Instead of directly searching the task-to-VM space, it dynamically selects the low-level heuristic that provides the most optimal result from a pool of scheduling algorithms.

The framework consists of a Python decision plane and a Java CloudSim plus simulation application which communicate via Inter Process Communication(IPC). The framework uses a multi-objective reward function that tanh squashes multiple execution metrics into a single, readable score between 0 and 1.

\subsection{Key Contributions}
The main contributions of this dissertation are:
\begin{itemize}
    \item \textbf{Decoupled Hyper-Heuristic Architecture:} We build a lightweight, streaming Python-CloudSim Plus framework that selects low-level heuristics dynamically without interrupting execution streams.
    \item \textbf{Non-Stationary Adaptation under Failure:} We demonstrate that Discounted UCB outperforms standard bandit variants during abrupt system failures, pivoting its policy from computationally intensive meta-heuristics to PEFT within 10 rounds after a 90\% VM capacity loss.
    \item \textbf{Multi-Objective Carbon-Aware Reward Model:} We design a batch-relative reward function using a hyperbolic tangent ($\tanh$) transformation to balance makespan, financial cost, carbon emissions, and resource utilization without scale distortion.
    \item \textbf{Trace-Driven Empirical Validation:} We evaluate our framework using real-world HPC2N and NASA workload traces across 10-seed multi-seed runs.
\end{itemize}

\subsection{Report Structure}
The rest of this report is organized as follows: Section~\ref{sec:rw} reviews related literature on cloud task scheduling, multi-armed bandits, and sustainable hyper-heuristics. Section~\ref{sec:Methodology} provides the theoretical foundation for our MAB formulations and reward functions. Section~\ref{sec:Design} presents the system design and pipeline architecture. Section~\ref{sec:Implementation} details the software implementation and IPC integration. Section~\ref{sec:Evaluation} presents experimental results, sensitivity analyses, and statistical evaluations. Finally, Section~\ref{sec:Conclusion} concludes the report and discusses directions for future work.